\documentclass[../main.tex]{subfiles}
\graphicspath{{\subfix{../image/}}}
\begin{document}
\chapter{Sard 定理}\label{chapter:ch06-Sard定理}

我们已经看到,正则值是我们最重要的工具之一,因为它们的水平集是光滑子流形。在本章,我们利用这一观察来证明光滑流形理论中两个最深刻的整体结果:每个紧光滑流形都可光滑嵌入某 Euclid 空间,以及任意光滑映射都可用满足附加正则性条件的映射来逼近。这些结果的核心是 Sard 定理,它刻画了光滑映射的临界值在何种意义下是稀有或可忽略的。

我们从零测集的定义开始,这是我们所说的``可忽略''的含义。然后陈述 Sard 定理并给出证明。我们应用它来证明紧流形的 Whitney 嵌入定理,随后应用于 Whitney 逼近定理,后者保证光滑映射可用光滑度任意高的映射来逼近,同时在紧集附近保持不变。我们以横截性理论的入门讨论结束本章,该理论表明横截于某子流形的映射在某种意义下是``一般的''。

\subfile{Chapter-06/section-01.tex}
\subfile{Chapter-06/section-02.tex}
\subfile{Chapter-06/section-03.tex}
\subfile{Chapter-06/section-04.tex}
\subfile{Chapter-06/section-05.tex}
\subfile{Chapter-06/section-06.tex}
\end{document}